\section{Internal structure of the electron}
\label{sec:internal-structure-of-the-electron}

In the brane framework, the electron is modeled as a stable, localized solitonic
excitation with toroidal topology. This section describes the geometric and
dynamical structure of this toroidal electron model, which serves as the
foundation for the numerical experiments described in
Section~\ref{experimental-setting}.

\subsection{Tubular world-tube description}
\label{subsec:tubular-electron}

The electron is realized as a toroidal soliton whose centerline traces a closed
curve in the brane. In the rest frame, this centerline lies in a fixed lateral
plane (say, the $X^1$--$X^2$ plane), forming a torus of major radius $R$ and
minor (tube) radius $r_0$. The torus axis defines the quantization axis for
spin.

We parametrize the centerline by arc length $z \in [0, 2\pi R)$ and introduce
tubular coordinates $(\rho, \theta)$ for the cross-section perpendicular to the
centerline at each point $z$, where $\rho$ is the radial distance from the
centerline and $\theta$ is the angular coordinate around the tube. The
displacement field $\xi(t,z,\rho,\theta)$ represents the amplitude of the
internal oscillatory mode circulating along the torus.

\subsection{Double-loop ridge profile (W\&vdM-inspired)}

Following the Wheeler--van der Mark picture, we model the internal photon-like
mode as a double-loop structure: two narrow ridges of high amplitude that spiral
around each other as one traverses the torus. Geometrically this may be viewed as a
twisted strip wrapped around the torus; when combined with a half-angle rotation of
the internal polarization frame, the resulting state exhibits Möbius-type (spinorial) $4\pi$
periodicity. The cross-sectional profile at
position $z$ along the centerline is given by

\begin{equation}
f(\rho,\theta;z) = G(\rho) \left[
e^{-(\theta-\alpha(z))^2/(2\sigma_\theta^2)}
+ e^{-(\theta-\alpha(z)-\pi)^2/(2\sigma_\theta^2)}
\right],
\label{eq:double-tunnel-profile}
\end{equation}

where:
\begin{itemize}
\item $G(\rho)$ is a radial envelope function (e.g., Gaussian) that confines the
  amplitude to a thin shell of radius $\rho \approx r_0$,
\item $\alpha(z) = \ell z / R$ controls the twist of the ridge pair as one moves
  along the torus (with $\ell$ the twist winding number),
\item $\sigma_\theta$ sets the angular width of each ridge.
\end{itemize}

With twist winding $\ell=1$, each ridge executes a single rotation as the torus
is traversed once, producing an intertwined $(2,1)$-type torus-knot pattern in
which the two ridges spiral around each other.


\subsection{Confinement by nonlinear self-guidance}
\label{subsec:electron-confinement}

A toroidal ``electron'' requires more than a prescribed geometric centerline: the internal
photon-like carrier must be \emph{guided} strongly enough to circulate on a ring of radius $R$
without rapidly dispersing or leaking away.
In this brane framework we do not introduce an external confinement potential or additional
``particle glue''. Instead, we hypothesize that confinement arises from the brane's constitutive
nonlinearity: the time-averaged stress generated by a large-amplitude, narrowband mode modifies the
local propagation characteristics of small fluctuations around the deformed state, producing a
self-induced waveguide.

It is useful to separate two statements. The \emph{shape} of the core (e.g.\ a toroidal bulge,
dominantly in the embedding-normal direction $X^4$) is a geometric property of the embedding
$\mathbf X$. By contrast, the effective guidance profile is encoded in a position-dependent
dispersion relation obtained from linearizing the dynamics around the background configuration; we
summarize this by an effective refractive-index profile $n_{\rm eff}(r)$ across the tube cross-section.
A finite maximal propagation speed of the medium does not by itself bend rays; bending requires a
transverse gradient of the effective wave speed.
Appendix~\ref{app:confinement-slice} presents a simple radial-slice model around the torus centerline,
yielding the order-of-magnitude requirement $\Delta n/n\sim a/R$ for a guided mode of transverse radius
$a$ on a ring of radius $R$.

For the double-ridge (double-winding) topology, coherence further constrains the cross-sectional
dispersion: the inner and outer ridges traverse different path lengths and therefore require different
local phase advance per arclength at fixed temporal frequency (Appendix~\ref{app:confinement-slice}).
In the remainder of this paper, the geometric ansatz provides an initialization target, while the
numerical experiments test whether the substrate dynamics admits a long-lived self-guided circulating
state consistent with these constraints.

\subsection{Compton-frequency oscillation}

The full displacement field for the electron soliton in its rest frame is

\begin{equation}
\xi(t,z,\rho,\theta) = A \, f(\rho,\theta;z) \cos(\omega_C t - m k_C z),
\label{eq:xi-electron-init}
\end{equation}

where:
\begin{itemize}
\item $A$ is the amplitude scale (of order the reduced Compton wavelength
  $\lambda_C = \hbar/(m_e c)$, as determined by the parameter matching in
  Section~\ref{amplitude-scale-calibration}),
\item $\omega_C = m_e c^2 / \hbar$ is the Compton angular frequency,
\item $k_C = 2\pi / \lambda_C$ is the Compton wavenumber,
\item $m$ is the longitudinal winding number (typically $m=2$, encoding that the
  internal Compton mode winds twice around the torus per revolution).
\end{itemize}

This configuration describes an internal standing wave pattern that circulates
at the microscopic wave speed $c$ along the toroidal world-tube.

\subsection{Dirac bilinears and effective four-current}

In the effective long-wavelength description, the electron is associated with a
Dirac spinor $\psi(x)$ whose bilinears encode the coarse-grained properties of
the soliton. For a rest-frame spin-up electron (spin aligned with the torus
axis, taken to be the $x^3$-direction), the Dirac bilinears are

\begin{equation}
j^\mu = \bar\psi \gamma^\mu \psi
= \begin{pmatrix} m_e c \\ 0 \\ 0 \\ 0 \end{pmatrix},
\qquad
s^i = \frac{1}{2} \bar\psi \gamma^i \gamma^5 \psi
= \begin{pmatrix} 0 \\ 0 \\ \hbar/2 \end{pmatrix}.
\label{eq:dirac-rest-bilinears}
\end{equation}

These bilinears serve as target constraints for initializing the toroidal ansatz
in numerical experiments: the axis of the torus, the direction of the internal
circulation of $\xi$, and the net tangential lateral velocity field in the tube
are chosen such that the resulting coarse-grained momentum and spin of the
soliton approximate these Dirac values. The subsequent evolution is always
governed by the full nonlinear brane equations; the Dirac constraints are used
only to fix and assess the macroscopic properties (charge, spin, effective
four-current) of the initial electron state.

\subsection{Spinorial $4\pi$ periodicity}

A key topological feature of the toroidal electron model is that the internal
frame (the orientation of the double-loop ridge pair) exhibits $4\pi$
periodicity under continuous rotation of the torus axis. This arises from the
twist of the ridge profile: as the torus is rigidly rotated by $2\pi$ about its
axis, the internal phase accumulated by the circulating mode does not return to
its initial value, but instead picks up a sign change. A full $4\pi$ rotation is
required to return the soliton to its original internal state.

In the language of geometric phase (Section~\ref{sec:brane-to-berry-to-em}), such a $4\pi$ periodicity would suggest that
loop transport around the defect probes a nontrivial internal frame.
A precise non-Abelian (Wilczek--Zee) formulation may become relevant if the excitation lives in an effectively degenerate
polarization subspace, but we do \emph{not} develop this machinery in the present step-by-step revision.
For now, $4\pi$ behaviour is treated as an empirical target signature for the discrete experiments, to be revisited once the
minimal $U(1)$ holonomy readout is numerically stable.

In the numerical programme, the tubular coordinates of Appendix~B are not merely
descriptive: they provide the scaffold for initializing a narrowband carrier
confined to a double-loop neighborhood and for sampling an internal polarization
frame along the loop. The spinorial signature is tested via a non-Abelian
Wilson-loop holonomy $W\in U(2)$ extracted from the transported degenerate
subspace (Section~\ref{sec:holonomy-wilson}), with the target $W\approx -I$ for
one circuit and $W^2\approx I$ for two.